\subsection{连续态射}

定理 1. 设 G 和 H 是在 R 或 \(Q_{p}\) 上的两个李群芽。设 f 是从 G 到 H 的连续态射。则 f 是解析的。

给 \(\mathrm{L}(\mathrm{G})\) 和 \(\mathrm{L}(\mathrm{H})\) 赋予定义其拓扑的范数，并使得对任意 \(\| [x,y]\| \leqslant \| x\| \| y\|\) 都有 \(x,y\)。存在 \(\mathrm{L}(\mathrm{G})\) 中以 0 为中心的开球 V，以及定义在 \(\phi\) 上的 \(\mathbf{G}\) 的指数映射 \(\mathbf{V}\)，满足：(1) \(\phi (\mathbf{V})\)\(e\) 是 G 中 e 的开邻域；(2) \(\phi\) 是从解析流形 V 到解析流形 \(\phi (\mathbf{V})\) 的同构；(3) 对所有 \(\phi (nx) = \phi (x)^n\) 和所有 \(x\in V\)，只要 \(n\in \mathbf{Z}\)，就有 \(nx\in V\)。对 H 类似地定义 W 和 \(\psi\)。必要时缩小 V，可设 \(f(\phi (\mathrm{V}))\subset \psi (\mathrm{W})\)。于是 \(g = \psi^{-1}\circ f\circ \phi\) 是从 V 到 W 的连续映射。

我们证明

\[
(x \in \mathrm{V}, \lambda \in \mathbf {Q} \text {   and   } \lambda x \in \mathrm{V}) \Rightarrow g (\lambda x) = \lambda g (x).\tag{1}
\]

可假设 \(\lambda \neq 0\)。写成 \(\lambda = \frac{r}{q}\)，其中 \(r, q\) 属于 \(\mathbf{Z} - \{0\}\)。令 \(y = \frac{r}{q} x\)。

若 \(\mathbf{K} = \mathbf{R}\)，记 \(z = \frac{x}{q} = \frac{y}{r} \in \mathrm{V}\)。于是 \(x = qz, y = rz\)、，从而

\[
g (x) = \psi^ {- 1} (f (\phi (q z))) = \psi^ {- 1} (f (\phi (z) ^ {q})) = \psi^ {- 1} (f (\phi (z)) ^ {q}).
\]

我们证明 \(\psi^{-1}(f(\phi(z))^q) = q\psi^{-1}(f(\phi(z))) = qg(z)\)。只需验证：若 \(u \in \psi(W)\) 且 \(u^q \in \psi(W)\)，则 \(\psi^{-1}(u^q) = q\psi^{-1}(u)\)；但若 \(u = \psi(v)\)、\(u^q = \psi(v')\)，则 \(v'/q \in W\)，且 \((\psi(v'/q))^q = u^q\)，从而 \(\psi(v'/q) = u = \psi(v)\)，所以 \(v' = qv\)。

同理，\(g(y) = rg(z)\)，从而 (1)。

若 \(\mathbf{K} = \mathbf{Q}_p\)，记 \(z = rx = qy \in \mathrm{V}\)，从而 \(g(z) = rg(x) = qg(y)\)，再次得到 (1)。

由于 \(\mathbf{Q}\) 在 K 中稠密，(1) 蕴含

\[
(x \in \mathrm{V}, \lambda \in \mathrm{K} \text {   and   } \lambda x \in \mathrm{V}) \Rightarrow g (\lambda x) = \lambda g (x).\tag{2}
\]

令 \(x \in L(G)\)，并令 \(\lambda, \lambda'\) 为 \(K^*\) 中的元素，使得 \(\lambda x \in V, \lambda' x \in V\)。则

\[
g (\lambda^ {\prime} x) = g \left(\frac {\lambda^ {\prime}}{\lambda} \lambda x\right) = \frac {\lambda^ {\prime}}{\lambda} g (\lambda x)
\]

由 (2)，可得 \(\frac{1}{\lambda} g(\lambda x) = \frac{1}{\lambda'} g(\lambda' x)\)\(h\)\(g\)。因此，g 到 \(\mathrm{L}(G)\) 的延拓 h 定义为：对于满足 \(h(x) = \frac{1}{\lambda} g(\lambda x)\) 的所有 \(\lambda\)，令 \(\lambda x \in \mathrm{V}\)\(h\)。显然 h 连续。我们证明

\[
(x \in \mathrm{L} (\mathrm{G}) \text {   and   } \lambda \in \mathrm{K}) \Rightarrow h (\lambda x) = \lambda h (x).\tag{3}
\]

令 \(\lambda' \in \mathbf{K}^*\)，使得 \(\lambda' x \in V\) 且 \(\lambda' \lambda x \in V\)。则

\[
h (\lambda x) = \frac {1}{\lambda^ {\prime}} g \left(\lambda^ {\prime} \lambda x\right) = \frac {1}{\lambda^ {\prime}} \lambda g \left(\lambda^ {\prime} x\right) = \lambda \frac {1}{\lambda^ {\prime}} g \left(\lambda^ {\prime} x\right) = \lambda h (x).
\]

令 \(x, y\) 属于 L(G)。则由 § 4，第 3 号，命题 4，

\[
\begin{array}{r l} h (x) + h (y) & = \lim _ {\lambda \in \mathbf {K} ^ {*},   \lambda \to 0} \lambda^ {- 1} \psi^ {- 1} (\psi (\lambda h (x)) \psi (\lambda h (y))) \\ & = \lim _ {\lambda \in \mathbf {K} ^ {*},   \lambda \to 0} \lambda^ {- 1} \psi^ {- 1} (\psi (h (\lambda x)) \psi (h (\lambda y))). \end{array}
\]

当 \(|\lambda|\) 充分小时，\(\lambda x \in V\) 且 \(\lambda y \in V\)，因而上式等于

\[
\begin{array}{r l} & \lim _ {\lambda \in \mathbf {K} ^ {*}, \lambda \to 0} \lambda^ {- 1} \psi^ {- 1} (f (\phi (\lambda x)) f (\phi (\lambda y))) \\ = & \lim _ {\lambda \in \mathbf {K} ^ {*}, \lambda \to 0} \lambda^ {- 1} (\psi^ {- 1} \circ f) (\phi (\lambda x) \phi (\lambda y)) \\ = & \lim _ {\lambda \in \mathbf {K} ^ {*}, \lambda \to 0} \lambda^ {- 1} g (\phi^ {- 1} (\phi (\lambda x) \phi (\lambda y))) \\ = & \lim _ {\lambda \in \mathbf {K} ^ {*}, \lambda \to 0} h (\lambda^ {- 1} \phi^ {- 1} (\phi (\lambda x) \phi (\lambda y))) \\ = & h (\lim _ {\lambda \in \mathbf {K} ^ {*}, \lambda \to 0} \lambda^ {- 1} \phi^ {- 1} (\phi (\lambda x) \phi (\lambda y))) \\ = & h (x + y). \end{array}
\]

因此 h 连续且线性，从而 \(h\)\(g = h|\mathrm{V}\)\(f\) 是解析的，继而 f 在 \(\phi(\mathrm{V})\)\(f\) 上解析，因而 f 是解析的（§ 1，第 10 号）。

推论 1. 设 G 是拓扑群。G 上至多存在一个与 G 的拓扑群结构相容的、在 R（分别地，\(Q_{p}\)）上的解析流形结构。

这立即由定理 1 得到。

定义 1. 若拓扑群 G 上存在一个与其拓扑相容的 R（分别地，p 进）李群结构，则称 G 为实（分别地，p 进）李群。

该结构于是唯一，我们可以谈论此类群的维数。若 G 和 H 是两个这样的群，则从 G 到 H 的每个连续态射都是解析的。

推论 2. 设 G 是拓扑群，V 是 e 的一个开邻域。假设 V 具有一个解析流形结构，使其成为实（分别地，p 进）李群芽。则 G 是实（分别地，p 进）李群。

设 \(g \in G\)。存在 G 中 e 的一个开邻域 \(V'\)\(e\)\(G\)，使得 \(V' \cup gV'g^{-1} \subset V\)。从 \(v \mapsto gvg^{-1}\) 到 V 的映射 \(V'\)\(V\) 是连续的，因而是从李群芽 \(V'\)\(V\) 到李群芽 V 的解析态射。因此只需应用 § 1，第 9 号，命题 18。

注记。(1) 若将 \(\mathbf{R}\)（分别地，\(\mathbf{Q}_p\)）替换为 \(\mathbf{C}\)，则定理 1 及其推论不再成立（练习 1）。

\begin{enumerate}
\def\labelenumi{(\arabic{enumi})}
\setcounter{enumi}{1}
\tightlist
\item
设 G 是拓扑群。可以证明†下列条件等价：(a) G 是有限维实李群；(b) G 是局部紧的，且存在 e 的一个邻域，其中不含 \{\} 以外的子群；(c) e 存在一个开邻域，同胚于空间 \(R^{n}\) 的开球。（一个容易得多的结果参见练习 6。）
\end{enumerate}
† 例如见 D. Montgomery 和 L. Zippin，《拓扑变换群》，纯粹与应用数学 Interscience 专论，第 1 号，Interscience 出版社，纽约，1955（尤其是第 169 和 184 页）。

命题 1. 设 \(G, G'\)\(f\)\(G\) 是拓扑群，f 是从 G 到 \(G'\) 的连续态射。假设下列三种情形之一成立：

\begin{enumerate}
\def\labelenumi{(\alph{enumi})}
\item
  G 是实李群，而 \(\mathbf{G}'\) 是 \(p\) 进李群；
\item
  G 是 p 进李群，而 \(G'\) 是实李群；
\item
  G 是 p 进李群，而 \(G'\) 是 \(p'\) 进李群，且 \(p \neq p'\)。
\end{enumerate}

则 \(f\)f 是局部常值的。

情形 (a)。令 \(G_0\)\(G\) 为 G 的单位分支。则 \(f(G_0)\) 是 \(G'\) 的连通子群，因而 \(f(G_0) = \{e\}\)，且 \(G_0\)\(G\) 在 G 中是开的。

情形 (b)。令 \(V'\) 为 \(G'\) 中 e 的一个邻域，使得包含于 \(G'\) 中的 \(V'\) 的每个子群都退化为 \(\{e\}\)（§4，第 2 号，定理 2 的推论 1）。存在 G 中 e 的一个邻域 V，使得 \(f(V) \subset V'\)。于是存在 G 的开子群 \(G_{1}\)，使得 \(G_{1} \subset V\)（§7，第 1 号，命题 1）。因此 \(f(G_{1}) = \{e\}\)。

情形 (c)。由 § 7，定理 4 及命题 8 的推论，存在  中 e 的一个邻域 \(\mathbf{V}'\)，使得对所有 \(e\)\(\mathbf{G}'\)\(x' \in \mathbf{V}' - \{e\}\)，当 n 趋于 \(x'^{p^n}\) 时 \(e\)\(n\)\(+\infty\) 不趋于 e。存在 G 中 e 的一个邻域 V，使得 \(V\)\(e\)\(G\)\(f(V) \subset V'\)。由 § 7，定理 4 及命题 9，存在 G 的开子群 \(G_1\)，使得 \(G\)\(G_1 \subset V\)，且对所有 \(x \in G_1\)，当 n 趋于 \(x^{p^n}\) 时 \(e\)\(n\)\(+\infty\) 趋于 e。因此 \(f(G_1) = \{e\}\)。

